% !TeX root = ../informe.tex
% ===========================================================================
%  5. Modelo de datos
%
%  El documento de alcance pide explícitamente el «modelo de datos» como
%  parte del entregable E1, y vale un 15 % de la rúbrica. Lo que se evalúa
%  no es solo que las tablas estén normalizadas, sino que se demuestre la
%  INDEPENDENCIA DE DATOS entre microservicios.
% ===========================================================================

\section{Modelo de datos}
\label{sec:modelo-datos}

\subsection{Estrategia de persistencia}
\label{sec:estrategia-datos}

\instruccion{Explica el reparto: qué base pertenece a qué servicio y qué
principio se sigue. Si se optó por una sola instancia de PostgreSQL con
varias bases en lugar de varias instancias, dilo aquí y justifícalo — es una
pregunta previsible en la defensa.}

Cada microservicio con estado es dueño de su propia base y nadie más la lee.
La regla no admite excepciones: si \id{ms-reservas} necesita el horario de una
cancha, se lo pregunta a \id{ms-canchas} por HTTP; no abre una conexión a
\id{canchas\_db}.

La decisión operativa es una \textbf{única instancia de PostgreSQL 16 con tres
bases independientes}, no tres instancias. Con tres instancias el aislamiento
sería más fuerte, pero el sistema entero tiene que levantarse con
\id{docker compose up} en la máquina de quien lo evalúa, y triplicar el motor
cuesta memoria y tiempo de arranque sin ganar nada demostrable en la rúbrica.

Lo que sí se conserva de un despliegue con instancias separadas es la parte que
importa: \textbf{un rol de base distinto por servicio, sin permisos cruzados}. Al
inicializar el volumen se revoca \id{CONNECT} a \id{PUBLIC} sobre cada base y se
concede solo a su dueño:

\begin{lstlisting}[language=SQL,caption={Aislamiento por credenciales, en \id{infra/postgres/init/01-bases-y-roles.sql}.},label={lst:roles}]
CREATE DATABASE usuarios_db OWNER usuarios_app;
REVOKE CONNECT ON DATABASE usuarios_db FROM PUBLIC;
GRANT  CONNECT ON DATABASE usuarios_db TO usuarios_app;
\end{lstlisting}

La diferencia es verificable: si alguien escribiera en \id{ms-reservas} una
consulta contra \id{canchas\_db}, no obtendría datos de otro dominio sino un
\id{FATAL: permission denied for database}. La independencia deja de ser una
convención que hay que recordar y pasa a ser algo que el motor impone.

\id{ms-reportes} no recibe rol ni credencial: no tiene base, agrega por REST.

\begin{table}[H]
  \centering
  \caption{Reparto de bases de datos.}
  \label{tab:bases}
  \begin{tabular}{L{3.2cm} L{3.4cm} L{6.6cm}}
    \toprule
    \textbf{Base} & \textbf{Servicio dueño} & \textbf{Contenido} \\
    \midrule
    \id{usuarios\_db} & \id{ms-usuarios} & Tabla \id{usuario}: credenciales (hash BCrypt), rol y estado de activación. \\
    \id{canchas\_db}  & \id{ms-canchas}  & Tablas \id{cancha} y \id{bloqueo\_mantenimiento}: catálogo, deporte, horario de atención y bloqueos. \\
    \id{reservas\_db} & \id{ms-reservas} & Tablas \id{reserva} y \id{configuracion}. Aloja la restricción de exclusión que implementa \rn{02}. \\
    \bottomrule
  \end{tabular}
\end{table}

\subsection{Diagrama entidad-relación}
\label{sec:der}

\instruccion{Un DER por base, o uno global con las fronteras entre bases
marcadas. Lo segundo comunica mejor la independencia, que es justo lo que se
evalúa.}

\begin{figure}[H]
  \centering
  \diagrama{modelo-datos}
  \caption{Modelo entidad-relación, con la frontera de cada base de datos.}
  \label{fig:der}
\end{figure}

\subsection{Diccionario de datos}
\label{sec:diccionario}

\instruccion{Una tabla por entidad. No hace falta repetir todos los tipos
SQL: basta el campo, su tipo lógico y qué significa. Los tipos exactos están
en los scripts del anexo.}

\begin{table}[H]
  \centering
  \caption{Entidad \id{reservas}.}
  \label{tab:dic-reservas}
  \begin{tabular}{L{3.4cm} L{3.2cm} L{6.6cm}}
    \toprule
    \textbf{Campo} & \textbf{Tipo} & \textbf{Descripción} \\
    \midrule
    \id{id}           & \id{bigint} identidad & Clave primaria. \\
    \id{cancha\_id}   & \id{bigint}    & Cancha reservada. \textbf{Sin clave foránea}: vive en otra base (\cref{sec:integridad}). \\
    \id{usuario\_id}  & \id{bigint}    & Dueño de la reserva. Sin clave foránea, por lo mismo. \\
    \id{fecha}        & \id{date}      & Día del bloque reservado. \\
    \id{hora\_inicio} & \id{time}      & Inicio del bloque, incluido. \\
    \id{hora\_fin}    & \id{time}      & Fin del bloque, excluido. \id{CHECK hora\_fin > hora\_inicio}. \\
    \id{estado}       & \id{varchar(20)} & \id{CONFIRMADA}, \id{CANCELADA} o \id{FINALIZADA}, con \id{CHECK}. En la tabla solo aparecen los dos primeros: \id{FINALIZADA} se deriva al leer. \\
    \id{creada\_en}   & \id{timestamp} & Instante de creación, con \id{DEFAULT now()}. \\
    \id{cancelada\_en} & \id{timestamp} & Nulo salvo que se haya cancelado. Alimenta el reporte de cancelaciones por período. \\
    \bottomrule
  \end{tabular}
\end{table}

\subsection{Integridad entre bases de datos}
\label{sec:integridad}

\instruccion{Es el apartado que más distingue un modelo de microservicios de
un monolito con varias tablas. Explica qué referencias cruzan la frontera
entre bases, por qué no pueden ser claves foráneas y dónde se valida
entonces la integridad.}

Dos columnas de \id{reserva} referencian entidades de otras bases:
\id{cancha\_id} apunta a \id{canchas\_db.cancha} y \id{usuario\_id} a
\id{usuarios\_db.usuario}. Ninguna de las dos es una clave foránea, y su
ausencia no es un descuido: una clave foránea es una restricción \emph{dentro
de una base}, y estas dos tablas están en bases distintas a las que la conexión
de \id{ms-reservas} ni siquiera puede acceder. No es que se haya decidido no
declararlas: es que no existe forma de declararlas.

La integridad se valida entonces donde sí puede validarse, antes de persistir:

\begin{itemize}
  \item \textbf{La cancha}: \id{BookingService} llama a \id{ms-canchas} por
        \id{CourtClientPort} y comprueba que existe, que está activa y que el
        bloque cae dentro de su horario de atención. Si no existe, la petición
        se rechaza con 404; si existe pero no admite la reserva, con 422.
  \item \textbf{El usuario}: no hace falta comprobarlo. El identificador llega
        en la cabecera \id{X-User-Id}, que el gateway escribe solo después de
        verificar la credencial contra \id{ms-usuarios}, y descarta cualquier
        \id{X-User-*} que venga del cliente. Un identificador inventado no puede
        llegar hasta aquí.
\end{itemize}

Si \id{ms-canchas} no responde, la reserva no se crea: se devuelve un 503. La
alternativa —persistir sin validar y reconciliar después— dejaría reservas sobre
canchas inexistentes o inactivas, que es exactamente el estado que la clave
foránea existía para impedir.

Inactivar un usuario o una cancha no toca las reservas ya hechas
(\textsf{FR-032}, \textsf{FR-047}): son de otro dominio, y esas bases no las
conocen.

\subsection{Restricciones que implementan reglas de negocio}
\label{sec:restricciones}

\instruccion{Aquí va la restricción de exclusión que impide el solapamiento.
Es el punto técnico más fuerte del modelo de datos: explica por qué está en
el motor y no en el código de la aplicación. El detalle de la regla como tal
va en la \cref{sec:solapamiento}; aquí solo el mecanismo de base de datos.}

La pieza principal es la restricción de exclusión de \id{reserva}, que impide
el solapamiento incluso bajo peticiones concurrentes. Su definición y el
razonamiento que la justifica están en la \cref{sec:solapamiento}; aquí interesa
el mecanismo: es un índice \textbf{GiST} sobre el par
$(\id{cancha\_id}, \id{tsrange(fecha + hora\_inicio, fecha + hora\_fin, '[)')})$,
con un predicado parcial que lo limita a las filas confirmadas. PostgreSQL
rechaza cualquier inserción cuyo rango se solape con uno ya indexado para la
misma cancha, y lo hace serializando en el índice, no en la aplicación. Requiere
la extensión \id{btree\_gist}, instalada al inicializar el volumen.

El resto de restricciones son \id{CHECK} que fijan las enumeraciones cerradas y
los rangos válidos, de modo que un valor inválido se rechace aunque el código lo
dejara pasar:

\begin{table}[H]
  \centering
  \caption{Restricciones e índices por tabla.}
  \label{tab:restricciones}
  \begin{tabular}{L{3.6cm} L{4.4cm} L{5.2cm}}
    \toprule
    \textbf{Tabla} & \textbf{Restricción o índice} & \textbf{Qué sostiene} \\
    \midrule
    \id{usuario} & \id{UNIQUE (username)}                       & \textsf{FR-002}: el duplicado es imposible, no solo improbable. \\
    \id{usuario} & \id{CHECK rol IN (...)}                      & Enumeración cerrada de roles. \\
    \id{cancha}  & \id{CHECK deporte IN (...)}                  & Enumeración cerrada de deportes. \\
    \id{cancha}  & \id{CHECK hora\_cierre > hora\_apertura}     & \rn{07}: un horario invertido no delimita bloques. \\
    \id{reserva} & \id{EXCLUDE USING gist (...)}                & \rn{02}, incluido el acceso concurrente. \\
    \id{reserva} & \id{CHECK estado IN (...)}                   & \rn{08}. \\
    \id{reserva} & \id{INDEX (usuario\_id, fecha)}              & Consulta «mis reservas». \\
    \id{reserva} & \id{INDEX (cancha\_id, fecha)}               & Consulta de disponibilidad. \\
    \bottomrule
  \end{tabular}
\end{table}

El esquema lo versiona Flyway y Hibernate arranca con
\id{ddl-auto: validate}: si el mapeo de una entidad y la tabla real divergen, el
servicio no arranca, en lugar de fallar más tarde con una columna inesperada.
